% Source: Solutions/Karan-Gangadharan-PS7-Q1-4.tex (LaTeX source). "Question 2(a)" renumbered to "Problem 7.2(a)".
\submittedby{Karan Gangadharan}

\begin{proof}
\textbf{(a)} Define an equivalence relation on $X=\R^2$ by
\[ (x_0,y_0)\sim(x_1,y_1)\iff x_0+y_0^2=x_1+y_1^2. \]
Let $X^*=X/{\sim}$ be the corresponding quotient space. We shall show that $X^*$ is
homeomorphic to $\R$.

Define a map $g\colon X\to\R$ by
\[ g(x,y)=x+y^2. \]
Since $g$ is a polynomial function, it is continuous. Also, $g$ is surjective, since for every
$t\in\R$, $g(t,0)=t$.

Now, define a map $f\colon\R\to X$ by $f(t)=(t,0)$. Clearly, $f$ is continuous and
\[ (g\circ f)(t)=g(t,0)=t. \]
Thus, $g\circ f=I_\R$. Therefore, by Problem 7.2(a), $g$ is a quotient map.

Let $q\colon X\to X^*$ be the natural quotient map defined by $q(x,y)=[(x,y)]$. Define
$\overline g\colon X^*\to\R$ by
\[ \overline g([(x,y)])=x+y^2. \]
We show that $\overline g$ is a homeomorphism.

\textit{(i) $\overline g$ is well-defined:}
Suppose $[(x_0,y_0)]=[(x_1,y_1)]$. Then, by the definition of the equivalence relation,
$x_0+y_0^2=x_1+y_1^2$. Hence,
$\overline g([(x_0,y_0)])=\overline g([(x_1,y_1)])$. Therefore, $\overline g$ is well-defined.

\textit{(ii) $\overline g$ is bijective:}
Suppose $\overline g([(x_0,y_0)])=\overline g([(x_1,y_1)])$. Then $x_0+y_0^2=x_1+y_1^2$. Thus,
$(x_0,y_0)\sim(x_1,y_1)$, which implies $[(x_0,y_0)]=[(x_1,y_1)]$. Therefore, $\overline g$ is
injective. Further, since $g$ is surjective, for every $t\in\R$, there exists $(x,y)\in X$ such
that $\overline g([(x,y)])=t$. Hence, $\overline g$ is surjective. Therefore, $\overline g$ is
bijective.

\textit{(iii) $\overline g$ is a homeomorphism:}
We have
\[ g=\overline g\circ q. \]
Since $g$ is continuous and $q$ is a quotient map, the quotient property of $q$ implies that
$\overline g$ is continuous.

Also, let $U$ be any open subset of $X^*$. Since $\overline g$ is injective, we have
\[ g^{-1}\bigl(\overline g(U)\bigr)=q^{-1}(U). \]
Since $q$ is a quotient map and $U$ is open in $X^*$, $q^{-1}(U)$ is open in $X$. Thus,
$g^{-1}(\overline g(U))$ is open in $X$. Since $g$ is a quotient map, $\overline g(U)$ is open
in $\R$. Therefore, $\overline g$ is an open map.

Hence, $\overline g$ is a continuous bijective open map, and consequently it is a homeomorphism.
Thus,
\[ \boxed{X^*\cong\R.} \]

\medskip
\textbf{(b)} Define an equivalence relation on $X=\R^2$ by
\[ (x_0,y_0)\sim(x_1,y_1)\iff x_0^2+y_0^2=x_1^2+y_1^2. \]
Let $X^*=X/{\sim}$ be the corresponding quotient space. We shall show that $X^*$ is
homeomorphic to $[0,\infty)$.

Define a map $g\colon X\to[0,\infty)$ by
\[ g(x,y)=x^2+y^2. \]
Since $g$ is a polynomial function, it is continuous. Also, $g$ is surjective, since for every
$t\in[0,\infty)$, $g(\sqrt t,0)=t$.

Now, define a map $f\colon[0,\infty)\to X$ by $f(t)=(\sqrt t,0)$. Clearly, $f$ is continuous
and
\[ (g\circ f)(t)=g(\sqrt t,0)=t. \]
Thus, $g\circ f=I_{[0,\infty)}$. Therefore, by Problem 7.2(a), $g$ is a quotient map.

Let $q\colon X\to X^*$ be the natural quotient map defined by $q(x,y)=[(x,y)]$. Define
$\overline g\colon X^*\to[0,\infty)$ by
\[ \overline g([(x,y)])=x^2+y^2. \]
We show that $\overline g$ is a homeomorphism.

\textit{(i) $\overline g$ is well-defined:}
Suppose $[(x_0,y_0)]=[(x_1,y_1)]$. Then, by the definition of the equivalence relation,
$x_0^2+y_0^2=x_1^2+y_1^2$. Hence,
$\overline g([(x_0,y_0)])=\overline g([(x_1,y_1)])$. Therefore, $\overline g$ is well-defined.

\textit{(ii) $\overline g$ is bijective:}
Suppose $\overline g([(x_0,y_0)])=\overline g([(x_1,y_1)])$. Then
$x_0^2+y_0^2=x_1^2+y_1^2$. Thus, $(x_0,y_0)\sim(x_1,y_1)$, which implies
$[(x_0,y_0)]=[(x_1,y_1)]$. Therefore, $\overline g$ is injective. Further, since $g$ is
surjective, for every $t\in[0,\infty)$, there exists $(x,y)\in X$ such that
$\overline g([(x,y)])=t$. Hence, $\overline g$ is surjective. Therefore, $\overline g$ is
bijective.

\textit{(iii) $\overline g$ is a homeomorphism:}
We have
\[ g=\overline g\circ q. \]
Since $g$ is continuous and $q$ is a quotient map, the quotient property of $q$ implies that
$\overline g$ is continuous.

Also, let $U$ be any open subset of $X^*$. Since $\overline g$ is injective, we have
\[ g^{-1}\bigl(\overline g(U)\bigr)=q^{-1}(U). \]
Since $q$ is a quotient map and $U$ is open in $X^*$, $q^{-1}(U)$ is open in $X$. Thus,
$g^{-1}(\overline g(U))$ is open in $X$. Since $g$ is a quotient map, $\overline g(U)$ is open
in $[0,\infty)$. Therefore, $\overline g$ is an open map.

Hence, $\overline g$ is a continuous bijective open map, and consequently it is a homeomorphism.
Thus,
\[ \boxed{X^*\cong[0,\infty).} \qedhere \]
\end{proof}
